% TEMPLATE-MILC-v3.tex - plantilla maestra UNICA de las guias MILC v3.
%
% La usan las cuatro familias de guias, cada una con su driver:
%   · Grados 6-11          -> scripts/build-guias-g11.py
%   · Bebras               -> scripts/build-guias-bebras.py
%   · Semillero            -> scripts/build-guias-semillero.py
%   · Territorio interior  -> scripts/build-guias-territorio-interior.py
%
% Antes habia tres copias casi identicas (~400 lineas iguales, 6 distintas):
% cada mejora habia que aplicarla tres veces, y en la practica no se hacia
% (el motor de recursos vivia solo en dos de ellas). Lo que varia entre
% familias esta parametrizado en 6 placeholders que cada driver llena
% (se nombran aqui SIN su sintaxis real: el builder reemplaza por texto plano,
% tambien dentro de los comentarios, y un valor multilinea rompe el preambulo):
%
%   HEADER_IZQ        encabezado izquierdo de pagina
%   HEADER_DER        encabezado derecho de pagina
%   PORTADA_ETIQUETA  rotulo sobre el numero grande (GRADO/MOMENTO/MODULO)
%   PORTADA_SERIE     linea de serie de la portada
%   PORTADA_ID        identificador de la guia en la portada
%   PORTADA_CIRCULO   el circulito de grado (°), o vacio si no aplica
%   PDF_TITULO        titulo en texto plano (sin \textbf ni comillas TeX) para
%                     los metadatos del PDF (/Title); lo produce lib_xelatex.texto_pdf
%
% Composicion (auditoria 2026-09-03): las bandas de estacion piden espacio
% (\Needspace*) y prohiben el salto justo despues (\nopagebreak), asi no quedan
% huerfanas al pie; las cajas largas (softbox, drawbox, infoband, puentebox) son
% `breakable`, asi una caja de media pagina ya no arrastra todo a la siguiente
% dejando la anterior al 40 %. Todo texto sobre mostaza o verde va en milcNegro
% (blanco sobre #E5B400 da 1,93:1 y sobre #58A923 2,95:1; ninguno pasa WCAG AA).
% El resto de claves las documenta content/guias/_SCHEMA.md. El script valida
% que no queden placeholders sin reemplazar antes de escribir el archivo final.

% Etiquetado PDF (latex-lab): /Lang, estructura y texto alternativo de las
% figuras, para lectores de pantalla. Debe ir ANTES de \documentclass.
\DocumentMetadata{lang=es-CO,tagging=on}
\documentclass[11pt,letterpaper]{article}
% headheight=24pt: la cabecera derecha lleva dos lineas (identidad de la guia
% y estacion actual). headsep se acorta en la misma medida para que el bloque
% de texto quede exactamente donde estaba con headheight=12pt.
\usepackage[margin=2.0cm,headheight=24pt,headsep=13pt]{geometry}
\usepackage[table]{xcolor}
\usepackage{fontspec}
\usepackage[spanish]{babel}
\usepackage{tikz}
% Librerías TikZ para los diagramas de las guías (bloque `recursos:`):
%  · babel  → babel en español vuelve activos caracteres como «>», y TikZ los usa
%             en sus opciones (>=stealth, ->). Sin esta librería, cualquier
%             diagrama con flechas revienta con «\language@active@arg> has an
%             extra }». Debe ir ANTES de que se dibuje nada.
%  · positioning        → sintaxis «below left=2pt and 3pt».
%  · decorations.pathreplacing → llaves ({brace}) para señalar rangos.
\usetikzlibrary{babel,positioning,decorations.pathreplacing}
\usepackage{graphicx}
% Circuitos: el trabajo de electrónica se hace hoy con micro:bit y sus sensores
% integrados (MakeCode, sin cableado), así que no se necesitan esquemáticos.
% Si algún día entran montajes con protoboard, basta descomentar esta línea:
% los diagramas .tex/.tikz declarados en `recursos:` ya se incrustan solos.
% \usepackage{circuitikz}
\usepackage[most]{tcolorbox}
\usepackage{tabularx}
\usepackage{array}
\usepackage{fancyhdr}
\usepackage{titlesec}
\usepackage[document]{ragged2e}  % bandera a la derecha en todo el cuerpo
\usepackage{enumitem}
\usepackage{microtype}
\usepackage{etoolbox}
\usepackage{needspace}
% hyperref al final del preambulo: metadatos del PDF (titulo, autor, idioma,
% familia), marcadores por estacion (\pdfbookmark) y enlaces sin recuadro.
\usepackage[hidelinks,
  pdftitle={Género y mandatos: quién decidió lo que se espera de ti},
  pdfauthor={PhD. Álvaro Cárdenas Orozco},
  pdfsubject={Territorio interior · Educación socioemocional · Grado 9° · Momento 2 · MILC v3},
  pdflang={es-CO},
  bookmarksopen=true,bookmarksnumbered=false]{hyperref}

% Helvetica Neue no trae la flecha U+2192, asi que cada «→» escrito literalmente
% en el YAML se compilaba a NADA, en silencio (462 flechas en 61 guias: rutas del
% tipo «paleta → tipografia → lockup» salian rotas). Mapeamos el caracter a su
% version matematica, que si tiene glifo. Asi el autor puede seguir escribiendo
% «→» comodamente en el YAML.
\usepackage{newunicodechar}
\newunicodechar{→}{\ensuremath{\rightarrow}}
% Mismo problema con los simbolos de marca y vinetas: el «✓» de «una palomita ✓»
% salia como cuadrito vacio en el PDF de todas las guias que lo usaban.
\usepackage{pifont}
\newunicodechar{✓}{\ding{51}}
\newunicodechar{✗}{\ding{55}}
\newunicodechar{✔}{\ding{52}}
\newunicodechar{•}{\textbullet}
\newunicodechar{≥}{\ensuremath{\geq}}
\newunicodechar{≤}{\ensuremath{\leq}}

\setmainfont{Helvetica Neue}
\setsansfont{Helvetica Neue}
\setlength{\parindent}{0pt}
\setlength{\parskip}{6pt}
% Sin particion de palabras: en 6.º-7.º una palabra cortada por guion al final
% de linea («Comuni-cacion») frena la lectura mucho mas que una linea corta.
\hyphenpenalty=10000
\exhyphenpenalty=10000
\pretolerance=2000
\tolerance=3000
\setlength{\emergencystretch}{3em}
\linespread{1.10}
\renewcommand{\arraystretch}{1.25}
\setlist{leftmargin=5mm,itemsep=1mm,topsep=1mm}
\newcolumntype{Y}{>{\RaggedRight\arraybackslash}X}

\definecolor{milcMagenta}{HTML}{E6007E}
\definecolor{milcVino}{HTML}{5A0038}
\definecolor{milcTurquesa}{HTML}{0093A5}
\definecolor{milcVerde}{HTML}{58A923}
\definecolor{milcMostaza}{HTML}{E5B400}
\definecolor{milcOcre}{HTML}{B86600}
\definecolor{milcNegro}{HTML}{241816}
\definecolor{milcGris}{HTML}{F7F7F4}
\definecolor{milcLinea}{HTML}{E8E2DD}
\definecolor{milcGradePrimary}{HTML}{0D9488}
\definecolor{milcGradeDark}{HTML}{115E59}
\definecolor{milcGradeSoft}{HTML}{CCFBF1}
\definecolor{milcGradeAccent}{HTML}{CFE7FF}
\definecolor{milcGradeText}{HTML}{FFFFFF}
\color{milcNegro}

\pagestyle{fancy}
\fancyhf{}
% Cabecera de dos lineas: identidad de la guia arriba y, a la derecha y en
% gris, la estacion en curso (\markright desde \milcestacion). Antes de la
% primera estacion la segunda linea va vacia; el \strut mantiene la altura
% para que la linea de arriba no baile entre paginas.
\lhead{\small Territorio interior · Educación socioemocional\\[-1pt]{\footnotesize\strut}}
\rhead{\small Grado 9° · Momento 2 · MILC v3\\[-1pt]{\footnotesize\strut\color{milcNegro!65}\rightmark}}
\cfoot{\small\thepage}
\renewcommand{\headrulewidth}{0pt}

% Sobre mostaza (#E5B400) y verde (#58A923) el texto va en milcNegro; sobre el
% resto de la paleta, en blanco. Se usa dentro de fonttitle/fontupper, donde un
% \color posterior manda sobre coltitle.
\newcommand{\milctintasobre}[1]{%
  \ifboolexpr{test{\ifstrequal{#1}{milcMostaza}} or test{\ifstrequal{#1}{milcVerde}}}%
    {\color{milcNegro}}{\color{white}}}

\newtcolorbox{titlebox}[1]{
  enhanced,colback=#1,colframe=#1,arc=7mm,boxrule=0pt,
  left=7mm,right=7mm,top=3mm,bottom=3mm,
  before skip=8pt,after skip=8pt,
  fontupper=\sffamily\bfseries\Large\color{white}
}

\newtcolorbox{softbox}[3]{
  enhanced,breakable,pad at break=3mm,
  colback=#2,colframe=#1,borderline west={4pt}{0pt}{#1},
  arc=2mm,boxrule=.4pt,left=4mm,right=4mm,top=3mm,bottom=3mm,
  before skip=6pt,after skip=8pt,title={#3},coltitle=white,
  colbacktitle=#1,fonttitle=\sffamily\bfseries\milctintasobre{#1},fontupper=\RaggedRight,
  attach boxed title to top left={xshift=3mm,yshift=-2mm},
  boxed title style={arc=2mm, boxrule=0pt}
}

\newtcolorbox{drawbox}[2]{
  enhanced,breakable,pad at break=4mm,
  colback=white,colframe=#1,arc=4mm,boxrule=1pt,
  left=5mm,right=5mm,top=4mm,bottom=4mm,before skip=8pt,after skip=8pt,
  title={#2},coltitle=white,colbacktitle=#1,fonttitle=\sffamily\bfseries\milctintasobre{#1},
  fontupper=\RaggedRight,
  attach boxed title to top left={xshift=4mm,yshift=-2mm},
  boxed title style={arc=3mm, boxrule=0pt}
}

\newtcolorbox{infoband}[1]{
  enhanced,breakable,pad at break=2mm,colback=#1!8,colframe=#1!50,arc=2mm,boxrule=.5pt,
  left=4mm,right=4mm,top=2mm,bottom=2mm,before skip=4pt,after skip=4pt,
  fontupper=\small\RaggedRight
}

% Puente narrativo entre secciones (contrato editorial Conéctate).
% Da continuidad: "Acabas de X. Ahora vas a Y porque Z."
% Visualmente: banda izquierda en color del grado, texto en itálica pequeña.
\newtcolorbox{puentebox}{
  enhanced,breakable,pad at break=2mm,colback=milcGradeSoft,colframe=milcGradeDark,
  borderline west={3pt}{0pt}{milcGradeDark},
  arc=0mm,boxrule=0pt,left=5mm,right=4mm,top=2.5mm,bottom=2.5mm,
  before skip=4pt,after skip=8pt,
  fontupper=\itshape\small\color{milcNegro}\RaggedRight
}
% La flecha va en modo matematico a proposito: Helvetica Neue NO tiene el glifo
% U+2192, asi que \textrightarrow se compilaba a NADA («Missing character: There
% is no → in font Helvetica Neue») y el puente salia sin flecha. En modo
% matematico la flecha viene de la fuente matematica y si se dibuja.
\newcommand{\puente}[1]{%
  \begin{puentebox}%
    \textcolor{milcGradeDark}{$\rightarrow$}\hspace{2mm}#1%
  \end{puentebox}%
}

\newcommand{\linead}[1]{\noindent\color{milcLinea}\rule{#1}{0.5pt}\color{milcNegro}}
\newcommand{\respuesta}[1]{\par\vspace{2mm}\foreach \n in {1,...,#1}{\noindent\color{milcLinea}\rule{\linewidth}{0.5pt}\par\vspace{5mm}}\color{milcNegro}}
\newcommand{\checkbox}{\textcolor{milcGradeDark}{\fbox{\phantom{x}}}\hspace{5pt}}

% ─── Listas aireadas para las guías socioemocionales ────────────────────
% Componer las verificaciones como «\checkbox texto\\» deja el texto que salta
% de línea debajo del cuadrito y apelmaza el bloque. Estos entornos alinean la
% continuación y airean los ítems. Los usa Territorio Interior; cualquier
% familia puede hacerlo.
\newlist{checklista}{itemize}{1}
\setlist[checklista]{
  label=\textcolor{milcGradeDark}{\fbox{\phantom{x}}},
  leftmargin=9mm, labelsep=3.2mm, itemsep=2.4mm,
  topsep=2.5mm, parsep=0pt, partopsep=0pt, before=\RaggedRight
}
\newlist{pasoslista}{enumerate}{1}
\setlist[pasoslista]{
  label=\textcolor{milcGradeDark}{\sffamily\bfseries\arabic*.},
  leftmargin=8mm, labelsep=3mm, itemsep=2.8mm,
  topsep=1mm, parsep=0pt, partopsep=0pt, before=\RaggedRight
}

\newcommand{\phasecard}[4]{
\begin{tcolorbox}[enhanced,colback=#3,colframe=#2,arc=3mm,boxrule=.5pt,height=3.05cm,valign=center,halign=center,left=1.5mm,right=1.5mm,top=2mm,bottom=2mm]
{\sffamily\bfseries\fontsize{22}{24}\selectfont\textcolor{#2}{#1}}\par
{\sffamily\bfseries\small #4}
\end{tcolorbox}
}

% ── Piezas del rediseno 2026 (legibilidad 6.º-11.º) ───────────────────
% Banda de estacion: reemplaza el titlebox macizo. Titulo a la izquierda,
% dato practico (tiempo/modalidad) a la derecha, en una sola linea.
% Nunca huerfana: pide 9 lineas libres (o salta de pagina rellenando con
% \vfill) y prohibe el salto justo despues de la banda. Nueve y no seis:
% la caja partible que sigue necesita ~80pt (titulo adosado, relleno y dos
% lineas) para arrancar en la misma pagina; con menos, tcolorbox la manda
% entera a la siguiente y la banda se queda sola. El penalty 10000 va
% en `after`, ANTES del espacio posterior: un `after skip` (aunque sea 0pt)
% es un \vskip tras la caja, y ese glue es un punto de corte legal que un
% \nopagebreak escrito despues de \end{tcolorbox} ya no cierra. Cada banda
% deja marcador PDF y actualiza la cabecera derecha.
\newcounter{milcestacion}
\newcommand{\milcestacion}[2]{%
\Needspace*{9\baselineskip}%
\stepcounter{milcestacion}%
\pdfbookmark[1]{#2}{milcestacion\themilcestacion}%
\markright{#2}%
\begin{tcolorbox}[enhanced,colback=#1,colframe=#1,arc=2mm,boxrule=0pt,
  left=5mm,right=5mm,top=2.6mm,bottom=2.6mm,before skip=11pt,
  after={\par\nopagebreak\vskip7pt}]
{\sffamily\bfseries\fontsize{15}{18}\selectfont\milctintasobre{#1}#2}
\end{tcolorbox}}

% Tarjeta de paso numerada: sustituye las tablas «Pilar / Aplicacion», que
% eran formato de consulta docente, no de estudiante. El numero va en un
% circulo sobre el borde para que la secuencia se lea de un vistazo.
\newcommand{\milcpaso}[3]{%
\Needspace{4\baselineskip}%
\begin{tcolorbox}[enhanced,colback=white,colframe=milcLinea,arc=1.5mm,boxrule=.9pt,
  left=14mm,right=4mm,top=3mm,bottom=3mm,before skip=4pt,after skip=4pt,
  fontupper=\RaggedRight,
  overlay={\node[anchor=north west,circle,fill=#2,text=white,inner sep=0pt,
    minimum size=7.4mm,font=\sffamily\bfseries\small]
    at ([xshift=3.6mm,yshift=-3.2mm]frame.north west) {#1};}]
#3
\end{tcolorbox}}

% Casilla marcable de tamano util con lapiz (la anterior era \fbox{\phantom{x}},
% de alto variable segun el texto que la seguia).
\newcommand{\milcmarca}{\framebox[3.8mm]{\rule{0pt}{3.4mm}}}

% Fila marcable de lista suelta (checks de la estacion 1).
\newcommand{\milccheckrow}[1]{%
\par\vspace{1.8mm}\noindent
\makebox[6.4mm][l]{\raisebox{-.55mm}{\textcolor{milcNegro}{\milcmarca}}}%
\parbox[t]{\dimexpr\linewidth-6.4mm\relax}{\RaggedRight #1}\par}

% Celda marcable dentro de la rubrica (tabla de autoevaluacion). Misma
% sangria francesa, pero pensada para vivir dentro de una columna X.
\newcommand{\milcopcion}[1]{%
\makebox[5.2mm][l]{\raisebox{-.55mm}{\textcolor{milcNegro}{\milcmarca}}}%
\parbox[t]{\dimexpr\linewidth-5.2mm\relax}{\RaggedRight #1}}

% ── Recursos visuales de la guía ──────────────────────────────────────
% Declarados en el bloque `recursos:` del YAML y emitidos por el builder.
% \guiaFigura   → imagen rasterizada o PDF (foto, carta celeste, montaje).
% \guiaDiagrama → fuente TikZ/circuitikz (.tex) incrustada como vector:
%                 es la vía para circuitos y esquemáticos editables.
% [#1] = texto alternativo (alt) · #2 = ruta relativa al .tex · #3 = pie
% El alt llega al PDF etiquetado (/Alt) y lo lee el lector de pantalla.
\newcommand{\guiaFigura}[3][]{%
\begin{center}
\includegraphics[width=0.88\linewidth,height=0.38\textheight,keepaspectratio,alt={#1}]{#2}\\[1.5mm]
{\footnotesize\itshape\textcolor{milcNegro!75}{#3}}
\end{center}
}

\newcommand{\guiaDiagrama}[2]{%
\begin{center}
\input{#1}\\[1.5mm]
{\footnotesize\itshape\textcolor{milcNegro!75}{#2}}
\end{center}
}

\begin{document}
\thispagestyle{empty}

% ── Portada ───────────────────────────────────────────────────────────
% Antes: un rectangulo de color a pagina completa. Se veia bien en pantalla,
% pero en la fotocopiadora del colegio se comia el toner y salia gris sucio.
% Ahora la banda ocupa el tercio superior y el resto va en blanco.
\begin{tikzpicture}[remember picture,overlay]
  \path[fill=milcGradePrimary]
    (current page.north west) rectangle ([yshift=-10.6cm]current page.north east);

  \node[anchor=north west,text=milcGradeText] at ([xshift=2.0cm,yshift=-1.55cm]current page.north west)
    {\sffamily\bfseries\fontsize{12}{14}\selectfont MOMENTO};
  \node[anchor=north west,text=milcGradeText,opacity=.85] at ([xshift=2.0cm,yshift=-2.35cm]current page.north west)
    {\sffamily\bfseries\fontsize{8.5}{10}\selectfont TERRITORIO INTERIOR · SOCIOEMOCIONAL};
  \node[anchor=north east,text=milcGradeText,opacity=.85,align=right] at ([xshift=-2.0cm,yshift=-1.55cm]current page.north east)
    {\sffamily\bfseries\fontsize{9}{12}\selectfont I.E. Sor María Juliana\\Cartago, Valle del Cauca};

  \node[anchor=west,text=milcGradeText] (gradonum) at ([xshift=1.85cm,yshift=-7.1cm]current page.north west)
    {\sffamily\bfseries\fontsize{112}{112}\selectfont 2};
  % El circulito de grado (°) solo aplica a las guias de grado («11°»); en
  % Bebras y Semillero el numero grande es un momento/modulo, no un grado.
  % Se posiciona relativo al nodo `gradonum`, no en coordenadas absolutas.
  

  \node[anchor=west,text=milcGradeText,text width=11.4cm,align=left] at ([xshift=1.5cm,yshift=0.5cm]gradonum.east)
    {
     {\sffamily\bfseries\fontsize{9}{11}\selectfont\MakeUppercase{Territorio interior · Socioemocional}}\\[3mm]
     {\sffamily\bfseries\fontsize{21}{25}\selectfont\setlength{\baselineskip}{27pt}Género\\[4pt]y mandatos\\[4pt]quién decidió lo que se espera}\\[3.5mm]
     {\sffamily\bfseries\fontsize{15}{18}\selectfont Grado 9° · Momento 2}
    };
\end{tikzpicture}

\vspace*{9.4cm}
\vfill

{\sffamily\bfseries\fontsize{8}{10}\selectfont\textcolor{milcGradeDark}{LO QUE TE LLEVAS HOY}}

\begin{tcolorbox}[enhanced,colback=white,colframe=milcNegro,arc=2mm,boxrule=1.6pt,
  left=5mm,right=5mm,top=4mm,bottom=4mm,before skip=3pt,after skip=12pt,
  fontupper=\RaggedRight\fontsize{11}{15.5}\selectfont]
Tu inventario de mandatos: las dos listas de lo que se espera de hombres y de mujeres en tu curso y en tu casa, con la cuenta de horas de trabajo no pagado de una semana real, y un mandato propio que vas a desobedecer con una acción concreta.
\end{tcolorbox}

\begin{tcolorbox}[enhanced,colback=milcGris,colframe=milcGris,arc=2mm,boxrule=0pt,
  left=5mm,right=5mm,top=3.5mm,bottom=3.5mm,after skip=12pt]
{\sffamily\bfseries\fontsize{8}{10}\selectfont\textcolor{milcNegro!55}{ESTA GUÍA ES DE}}\\[3mm]
\begin{tabularx}{\linewidth}{X p{3.2cm} p{3.2cm}}
\linead{\linewidth} & \linead{3.0cm} & \linead{3.0cm}\\[-1mm]
{\fontsize{8.5}{10}\selectfont\textcolor{milcNegro!55}{Nombre completo}} &
{\fontsize{8.5}{10}\selectfont\textcolor{milcNegro!55}{Grupo}} &
{\fontsize{8.5}{10}\selectfont\textcolor{milcNegro!55}{Fecha}}\\
\end{tabularx}
\end{tcolorbox}

\vfill

\noindent\textcolor{milcLinea}{\rule{\linewidth}{1pt}}\\[2mm]
{\sffamily\bfseries\fontsize{9.5}{12}\selectfont Autor: Dr. Álvaro Cárdenas Orozco}\\[1mm]
{\sffamily\fontsize{8.5}{11}\selectfont\textcolor{milcNegro!55}{Metodología MILC · Apertura ancestral · Género y mandatos · Triángulo Dussel-Estoicismo-Floridi}}

\newpage

\section*{Apertura: Género y mandatos: quién decidió lo que se espera de ti}

\begin{softbox}{milcGradeDark}{milcGradeSoft}{Saber ancestral}
En la finca cafetera ---la economía que hizo este territorio--- \textbf{el trabajo estuvo repartido por género, y en buena medida lo sigue estando}. Los estudios sobre el campo colombiano lo describen así: \emph{los hombres se ocupan preferentemente de las labores agrícolas, mientras las mujeres atienden las tareas de la casa y del entorno cercano a la vivienda, además del cuidado de los otros miembros de la familia}. \textbf{Y hay cifras que muestran cuánto de eso sigue vivo:} el \textbf{92\,\%} de las mujeres se dedica a la \textbf{preparación de alimentos}, mientras el \textbf{83\,\%} de los hombres se ocupa de \textbf{aplicar agroquímicos}. \textbf{Las mujeres son el 30\,\% de los caficultores del país} ---163.046 personas--- y su participación es fuerte, sobre todo en cosecha. \textbf{Ahora el dato que hay que subrayar,} porque es el corazón de esta guía: \emph{esa alta participación de las mujeres en la producción \textbf{no vino acompañada de ningún cambio en el reparto del trabajo doméstico}}, que siguió siendo responsabilidad suya. \textbf{Es decir: se sumó una jornada, no se repartió la otra.} \emph{Nadie firmó eso. No hay una ley que lo ordene. Y sin embargo se cumple.}
\end{softbox}

\begin{softbox}{milcTurquesa}{milcGris}{Saber contemporáneo}
¿Y cuándo se aprende lo que «corresponde» a cada uno? \textbf{Mucho antes de lo que uno creería.} \textbf{Lin Bian}, \textbf{Sarah-Jane Leslie} y \textbf{Andrei Cimpian} estudiaron a niñas y niños de \textbf{5, 6 y 7 años} y les preguntaron, con historias, quién era «\emph{muy, muy inteligente}» ---la versión infantil de lo que los adultos llaman brillantez---. \textbf{A los 5 años no había diferencia.} \textbf{A los 6, sí:} las niñas ya asociaban esa cualidad \textbf{más a los hombres y a los niños} que a las mujeres y a las niñas. \textbf{Y el efecto se veía en lo que querían hacer:} las niñas se mostraron \textbf{significativamente menos interesadas} que los niños en un juego presentado como «\emph{para niños muy inteligentes}». \textbf{Aquí viene el detalle que lo vuelve demoledor:} \emph{cuando el mismo juego se presentaba como «para niños que se esfuerzan mucho», \textbf{la diferencia desaparecía}} (Bian, Leslie y Cimpian, 2017). \emph{Es decir: no era que las niñas quisieran menos el juego. Era que no se creían con derecho al que pedía \textbf{ser brillante}.} \textbf{En un año. Sin que nadie les dijera nada explícitamente.}
\end{softbox}

\begin{softbox}{milcMostaza}{milcGris}{Pregunta-puente}
\textit{Nadie firmó que la cocina fuera de las mujeres, y sin embargo el 92\,\% la atiende. \textbf{¿Qué cosa haces o dejas de hacer «porque sí»} \ldots \textbf{que si fueras del otro género probablemente harías al revés}?}
\end{softbox}

\begin{softbox}{milcVerde}{milcGris}{Saber y saber hacer}
Al terminar podrás: (1) \textbf{reconocer} un mandato de género por su forma ---no se anuncia, se da por obvio y se castiga cuando se rompe---; (2) \textbf{explicar} que estas expectativas se instalan muy temprano y que afectan lo que alguien se cree con derecho a intentar; (3) \textbf{contar} las horas reales de trabajo no pagado en tu casa durante una semana; (4) \textbf{escoger} un mandato que te apriete a ti y desobedecerlo con una acción concreta, no con un discurso.
\end{softbox}



\small
\begin{tabularx}{\textwidth}{>{\bfseries}p{3.0cm}Y}
\rowcolor{milcGradePrimary!12}Autor & Dr. Álvaro Cárdenas Orozco\\
DBA & Comprende los fenómenos que afectan a su territorio, regula sus emociones ante ellos y acompaña a otros con empatía (transversal 6°--11°, articulado con la política «Escuela, territorio de vida» del MEN, Resolución 006519 de 2025, y la Ley 1620 de 2013)\\
Referentes & Primera ayuda psicológica (OMS, 2011/2012) · Directrices IASC (2007) · USGS y Servicio Geológico Colombiano · Pueblo nasa y la minga (CRIC) · Dussel · Estoicismo · Floridi\\
Producto final & Tu inventario de mandatos: las dos listas de lo que se espera de hombres y de mujeres en tu curso y en tu casa, con la cuenta de horas de trabajo no pagado de una semana real, y un mandato propio que vas a desobedecer con una acción concreta.\\
\end{tabularx}
\normalsize

\puente{En el momento 1 aprendiste a mirar desde el lugar de otro. \textbf{Hoy miras algo que ya estaba puesto antes de que llegaras:} lo que se espera de ti por haber nacido con un cuerpo o con otro.}

\milcestacion{milcGradeDark}{Ruta MILC: así vamos a aprender}
\begin{tabularx}{\textwidth}{>{\centering\arraybackslash}X>{\centering\arraybackslash}X>{\centering\arraybackslash}X>{\centering\arraybackslash}X}
\phasecard{1}{milcVerde}{milcGris}{Escucha\\[-1mm]\small Observo, escucho y pregunto.} &
\phasecard{2}{milcTurquesa}{milcGris}{Sistematización\\[-1mm]\small Nombro lo que se espera.} &
\phasecard{3}{milcMagenta}{milcGris}{Praxis\\[-1mm]\small Construyo y aplico.} &
\phasecard{4}{milcVino}{milcGris}{Evaluación\\[-1mm]\small Reflexión liberadora.}
\end{tabularx}


\puente{Primero, las dos listas ---lo que se espera de ellos y de ellas--- escritas sin adornar. \emph{Todos las conocen; casi nunca se escriben.}}

\milcestacion{milcVerde}{Estación 1 de 3 · Escucha}

\begin{softbox}{milcVerde}{milcGris}{Escena para escuchar}
\textbf{Actividad 1 · IDENTIFICA --- Las dos listas.} \textbf{Nadie va a leer esta página.} \textbf{Primera parte.} Haz \textbf{dos listas}: \emph{lo que se espera de un hombre} y \emph{lo que se espera de una mujer} en tu curso, en tu casa y en tu barrio. \textbf{Escríbelo tal como se dice}, aunque suene feo: «que no llore», «que ayude en la cocina», «que se haga respetar», «que sea juicioso», «que no ande sola de noche», «que no se deje». \textbf{Segunda parte.} Al frente de cada cosa marca \textbf{qué pasa si alguien no lo cumple}: ¿le dicen algo? ¿se ríen? ¿le ponen un apodo? ¿no pasa nada? \emph{Ese castigo es la prueba de que era un mandato y no una costumbre.} \textbf{Tercera parte.} En tu casa, \textbf{¿quién cocina, quién lava, quién plancha, quién cuida a los menores o a los mayores, quién saca la basura, quién arregla lo que se daña?} \textbf{Y la última:} ¿hay algo que te gustaría hacer y no haces \textbf{por lo que dirían}?
\end{softbox}

{\sffamily\bfseries\fontsize{8}{10}\selectfont\textcolor{milcNegro!55}{MARCA CUANDO LO HAYAS HECHO}}
\milccheckrow{Escribiste las dos listas tal como se dicen las cosas, sin suavizarlas.}
\milccheckrow{Anotaste \textbf{qué pasa cuando alguien no cumple} cada una.}
\milccheckrow{Escribiste quién hace cada oficio en tu casa.}

\begin{infoband}{milcVerde}


\textbf{En tu cuaderno (Actividad 1 · IDENTIFICA).} Título: «Actividad 1. Las dos listas». \textbf{Formato:} las dos listas con el castigo al frente + el reparto de oficios de tu casa + lo que no haces por lo que dirían. \textbf{Extensión:} media página. \textbf{Sabes que terminaste cuando} identificaste \textbf{al menos un mandato que te aprieta a ti}. Esta página es tuya: \textbf{nadie más tiene que leerla si tú no quieres}.
\end{infoband}

\puente{Ya las tienes. Ahora, cómo funciona un mandato, a qué edad llega y qué trabajo es el que nunca se cuenta.}

\milcestacion{milcTurquesa}{Fase 2 · Sistematización: entender el mandato}

\begin{softbox}{milcTurquesa}{milcGris}{Cuatro claves sobre los mandatos}
Cuatro claves para ver algo que funciona precisamente porque no se ve.
\end{softbox}

\milcpaso{1}{milcTurquesa}{\textbf{Un mandato no se anuncia: se da por obvio, y se cobra cuando se rompe.} \emph{Nadie te dijo nunca «los hombres no lloran»} como quien lee un reglamento: te lo enseñaron con una risa, con un silencio incómodo, con un apodo. \textbf{Por eso la prueba para detectar un mandato no es buscar quién lo dijo ---nadie lo dijo--- sino mirar \emph{qué pasa cuando alguien no lo cumple}.} \emph{Si no pasa nada, era una costumbre. Si hay burla, corrección o castigo, era un mandato.} \textbf{Y fíjate en el parecido con lo que ya viste:} \emph{funciona igual que la cerca del momento 8 del grado 8 ---nadie la construyó, nadie recuerda de dónde salió, y todos la respetan---.}}
\milcpaso{2}{milcTurquesa}{\textbf{Llega antes de los siete años.} El estudio de Bian, Leslie y Cimpian encontró que a los \textbf{5 años} niñas y niños no diferenciaban, y que \textbf{a los 6 las niñas ya asociaban «muy, muy inteligente» más a los niños}; y que se apartaban del juego «para muy inteligentes» \textbf{aunque no del juego «para los que se esfuerzan»} (Bian, Leslie y Cimpian, 2017). \emph{Detente en esa distinción, porque es lo más importante de la guía:} \textbf{el mandato no les quitó las ganas de jugar. Les quitó la idea de tener derecho a lo que exige ser brillante.} \emph{Y esa es la forma en que un mandato hace su verdadero daño:} \textbf{no te prohíbe nada ---te convence de que eso no es para ti, y entonces ni lo intentas---.}}
\milcpaso{3}{milcTurquesa}{\textbf{El trabajo que no se cuenta.} En el campo colombiano las mujeres entraron a la producción ---son el \textbf{30\,\%} de los caficultores--- \textbf{y el trabajo de la casa siguió siendo suyo}: la participación creció, el reparto doméstico no cambió. \emph{A eso se le llama, con razón, una \textbf{doble jornada}.} \textbf{Y hay una razón por la que es tan difícil discutirlo:} \emph{el trabajo doméstico y de cuidado \textbf{no se paga, no aparece en ninguna cuenta y no se nombra como trabajo}} ---se llama «ayudar»---. \textbf{Por eso este momento pide contar horas:} \emph{lo que no se mide parece que no existe, y una cifra hace lo que no logra ninguna discusión}.}
\milcpaso{4}{milcTurquesa}{\textbf{También aprieta a ellos, aunque de otra manera.} \emph{Sería un error salir de aquí pensando que el mandato solo pesa sobre las mujeres.} \textbf{A los hombres el mandato les cobra otras cosas:} no pedir ayuda, no llorar, no decir que algo les duele, «no dejarse», responder con rabia porque es la única emoción permitida. \emph{Y eso tiene consecuencias que ya trabajaste en otros momentos:} \textbf{el que no puede pedir ayuda es el que no la pide cuando la necesita} ---y en el grado 10 vas a ver lo que eso cuesta---. \textbf{El mandato es el mismo mecanismo con dos caras:} \emph{a unas les quita el derecho a intentar; a otros, el derecho a necesitar.}}

\begin{drawbox}{milcTurquesa}{Cómo se reconoce un mandato}
\begin{pasoslista}
\item \textbf{Nadie lo escribió,} y sin embargo todos lo saben.
\item \textbf{Se presenta como naturaleza:} «es que las mujeres son más\ldots», «es que los hombres son así».
\item \textbf{Hay castigo si se rompe:} burla, apodo, corrección, silencio. \emph{Sin castigo no es mandato.}
\item \textbf{Reparte trabajo o reparte permiso:} \emph{quién hace qué, y quién puede aspirar a qué.}
\item \textbf{Prueba final:} \emph{cámbiale el género a la frase. Si suena absurda al revés, era un mandato.}
\end{pasoslista}
\end{drawbox}

\begin{softbox}{milcMostaza}{milcGris}{Errores comunes y su corrección}
\textbf{«Eso ya cambió»:} el 92 \% y el 83 \% son de hoy, no de 1950. \textbf{«En mi casa no pasa»:} cuenta las horas antes de afirmarlo. \textbf{Creer que solo aprieta a las mujeres:} a ellos les cobra el derecho a necesitar. \textbf{Confundir mandato con gusto:} si a alguien le gusta cocinar, perfecto; la pregunta es si podría no hacerlo. \textbf{Discutirlo con discurso:} las horas contadas convencen más. \textbf{Escoger un mandato ajeno para desobedecer:} el ejercicio es sobre el tuyo. \textbf{«Es la costumbre»:} la costumbre no castiga; el mandato sí.
\end{softbox}

\begin{infoband}{milcTurquesa}


\textbf{En tu cuaderno (Actividad 2 · EXPLICA).} Título: «Actividad 2. La prueba del cambio de género». \textbf{Formato:} toma \textbf{cinco} frases de tus listas y escríbelas con el género cambiado; anota cuáles suenan absurdas al revés y cuáles no. \textbf{Extensión:} media página. \textbf{Sabes que terminaste cuando} puedes explicar por qué \textbf{un mandato no prohíbe: convence de que algo no es para uno}.
\end{infoband}

\puente{Ya sabes reconocerlos. Es hora de lo único que mueve algo: \textbf{contar horas} y \textbf{desobedecer uno} ---el tuyo---.}

\milcestacion{milcMagenta}{Fase 3 · Praxis: desobedecer uno}

\begin{softbox}{milcMagenta}{milcGris}{Cómo se desobedece un mandato sin discurso}
Dos cosas concretas: un número que nadie puede discutir y una acción que solo depende de ti.
\end{softbox}

\milcpaso{1}{milcMagenta}{\textbf{Las dos listas, terminadas (7 min).} Complétalas con el castigo al frente de cada mandato. \emph{Compara con dos compañeros ---uno de cada género si se puede---:} \textbf{van a descubrir que las listas se parecen muchísimo entre casas distintas}, y eso es exactamente lo que prueba que no son gustos personales.}
\milcpaso{2}{milcMagenta}{\textbf{La cuenta de horas (una semana, 3 min al día).} Durante siete días anota \textbf{quién hizo cada oficio en tu casa y cuánto tardó}: cocinar, lavar loza, lavar ropa, planchar, barrer, mercar, cuidar a alguien, ayudar con tareas a un hermano. \textbf{Al final, suma por persona.} \emph{No hace falta que sea exacto ---basta con que sea honesto---.} \textbf{Ese total es el producto más contundente de esta guía:} \emph{una cifra de tu propia casa.}}
\milcpaso{3}{milcMagenta}{\textbf{El mandato que me toca (8 min).} De tus listas, escoge \textbf{uno que te apriete a ti} ---no uno que le apriete a otro---. \emph{Puede ser pequeño:} no llorar delante de nadie, no cocinar nunca, no pedir ayuda en matemáticas, no bailar, no jugar cierto deporte, no salir sola, no decir que algo te dio miedo. \textbf{Escribe qué te ha costado ese mandato en concreto.}}
\milcpaso{4}{milcMagenta}{\textbf{La acción concreta (10 min).} Escribe \textbf{una sola acción} para desobedecerlo esta semana, con \textbf{día y hora}. \emph{Y que sea acción, no opinión:} \textbf{aprender a hacer algo que «no te corresponde», pedir la ayuda que no pides, asumir un oficio de la casa que nunca has hecho, inscribirte en algo}. \textbf{Después anota qué pasó y ---esto es lo interesante--- \emph{quién dijo algo}.} \emph{Si nadie dijo nada, también anótalo: es el hallazgo más frecuente y el más liberador.}}

\begin{drawbox}{milcMagenta}{Antes de cerrar tu inventario}
\begin{checklista}
\item Tus dos listas tienen el \textbf{castigo} anotado al frente de cada mandato.
\item Pasaste cinco frases por la \textbf{prueba del cambio de género}.
\item Contaste las horas \textbf{siete días} y sumaste por persona.
\item El mandato que escogiste \textbf{te aprieta a ti}, no a otro.
\item Tu acción tiene \textbf{día y hora}, y es acción y no opinión.
\item Anotaste qué pasó y \textbf{quién dijo algo} ---incluido si nadie dijo nada---.
\end{checklista}
\end{drawbox}

\begin{softbox}{milcMagenta}{milcGris}{Plantilla de trabajo}
\textbf{La prueba del cambio de género.} \emph{«\_\_\_\_» → al revés suena: «\_\_\_\_». ¿Absurda? Sí / No.} \\[1mm]
\textbf{La cuenta.} \emph{Persona · oficio · minutos · total de la semana.} \\[1mm]
\textbf{Mi desobediencia.} \emph{«El mandato que me aprieta es \_\_\_\_. Voy a \_\_\_\_ el \_\_\_\_ a las \_\_\_\_.»}
\end{softbox}

\begin{infoband}{milcMagenta}


\textbf{En tu cuaderno (Actividad 3 · CREA).} Título: «Actividad 3. Mi inventario de mandatos». \textbf{Formato:} las dos listas con castigos + la prueba del cambio de género + la tabla de horas sumada por persona + el mandato propio y la acción con su resultado. \textbf{Extensión:} una página más la tabla. \textbf{Sabes que terminaste cuando} tienes \textbf{el total de horas por persona} de tu casa.
\end{infoband}

\puente{El producto tiene un número adentro. \emph{Las horas contadas convencen a gente que ningún argumento convence.}}

\milcestacion{milcMostaza}{Producto de la sesión}
\textbf{Inventario de mandatos: las dos listas con su castigo, la cuenta de horas de una semana y un mandato desobedecido}.
\begin{softbox}{milcMostaza}{milcGris}{Criterios de calidad}
(1) Las listas registran los mandatos tal como se enuncian, con el castigo que sigue a incumplirlos. (2) La prueba del cambio de género está aplicada a cinco frases con su resultado. (3) La cuenta de horas cubre siete días y está sumada por persona. (4) El mandato elegido para desobedecer aprieta a quien hace el ejercicio, no a un tercero. (5) La desobediencia es una acción fechada, con registro de lo que ocurrió después.
\end{softbox}

\puente{Antes de cerrar, revisa lo que hace fracasar este ejercicio: si el mandato que escogiste desobedecer es uno que le aprieta a otro y no a ti, no cuenta.}

\milcestacion{milcVino}{Evaluación liberadora · cinco dimensiones}

\begin{softbox}{milcVino}{milcGris}{Sentido de la evaluación}
Evaluar no es solo revisar una nota. Es reconocer qué aprendí, qué cambió en mí, qué emociones manejé, cómo me relaciono con la ciudad y la comunidad, y qué le dejaré a quien viene detrás.
\end{softbox}

\textbf{Mi reflexión liberadora en cinco dimensiones.} Completa cada frase iniciada:

\small
\begin{tabularx}{\textwidth}{>{\bfseries\RaggedRight\arraybackslash}p{3.4cm}Y>{\RaggedRight\arraybackslash}p{4.6cm}}
\rowcolor{milcVino}\color{white}Dimensión & \color{white}\bfseries Mi respuesta (frame starter) & \color{white}\bfseries Algo que descubrí\\
1. Personal & Lo que más cambió en mí fue \linead{6.5cm} & \linead{4.2cm}\\
2. Emocional & La emoción más fuerte fue \linead{2.5cm} y la regulé \linead{3cm} & \linead{4.2cm}\\
3. Ciudadana & En democracia esto importa porque \linead{6cm} & \linead{4.2cm}\\
4. Local (Cartago) & En mi barrio sería útil para \linead{6.5cm} & \linead{4.2cm}\\
5. Intergeneracional & A mi abuela/abuelo le diría \linead{6.5cm} & \linead{4.2cm}\\
\end{tabularx}
\normalsize

\begin{infoband}{milcVino}
\textbf{Para la próxima sustentación.} Vuelve a esta tabla en seis meses. Verás que las respuestas cambian — no porque lo aprendido cambie, sino porque tú cambias. La evaluación liberadora es una conversación contigo a través del tiempo.
\end{infoband}


\puente{Cerramos con tres voces: el rostro que reclama un derecho, lo que sí está en tu poder, y el mundo compartido donde también se aprende qué se espera.}

\milcestacion{milcGradeDark}{Triángulo de pensamiento · cierre filosófico}

\begin{softbox}{milcVino}{milcGris}{Dussel · lente del nosotros}
\textit{«El otro se revela realmente como otro\ldots como el pobre, el oprimido; el que, a la vera del camino, fuera del sistema, muestra su rostro sufriente y sin embargo desafiante: «¡Tengo hambre!, ¡tengo derecho a comer!»»}

{\footnotesize\itshape\color{milcNegro!70}--- Enrique Dussel · Filosofía de la liberación (1977)}

Dussel llama la atención sobre \textbf{quién queda fuera del sistema}, y un reparto de mandatos es un sistema aunque nadie lo haya escrito. \emph{Fíjate en el trabajo doméstico:} \textbf{sostiene todo lo demás ---sin comida y sin ropa limpia no hay finca, ni colegio, ni oficina--- y sin embargo no aparece en ninguna cuenta}. \emph{Es trabajo que está afuera: no se paga, no se nombra, no se agradece.} \textbf{Y la última palabra de la cita vuelve a ser la clave:} \emph{quien pide que se reparta no está pidiendo un favor} ---no está pidiendo que «le ayuden» en su casa: está diciendo que ese trabajo \textbf{no era suyo solo}---.

\textbf{En mi casa, ¿de quién es el trabajo que sostiene todo y no se cuenta como trabajo?}
\end{softbox}
\linead{\linewidth}\\[1mm]
\linead{\linewidth}

\begin{softbox}{milcMostaza}{milcGris}{Estoicismo · Epicteto · Enquiridión, 1 (c. 125 d.C.; trad. desde E. Carter) · cuidado interior}
\textit{«Algunas cosas están en nuestro poder y otras no. Están en nuestro poder la opinión, el impulso, el deseo, el rechazo; en una palabra, todo aquello que es obra nuestra.»}

{\footnotesize\itshape\color{milcNegro!70}--- Epicteto · Enquiridión, 1 (c. 125 d.C.; trad. desde E. Carter)}

Un mandato se siente como si viniera de afuera y fuera inevitable. \emph{Y en parte lo es:} \textbf{no está en tu poder cambiar lo que la gente espera de ti}, ni evitar que alguien comente. \textbf{Pero fíjate en lo que sí lo está, según Epicteto:} \emph{la opinión ---qué tan cierto te parece el mandato--- y el impulso ---lo que haces---.} \textbf{Y hay una segunda lectura, más incómoda y más útil:} \emph{si el impulso es obra tuya, entonces también es obra tuya \textbf{el castigo que aplicas} cuando alguien rompe un mandato} ---la risita, el apodo, el silencio---. \textbf{Nadie sostiene un mandato solo. Lo sostienen los que corrigen.}

\textbf{¿A quién he corregido, aunque sea con una risa, por no comportarse «como corresponde»?}
\end{softbox}
\linead{\linewidth}\\[1mm]
\linead{\linewidth}

\begin{softbox}{milcTurquesa}{milcGris}{Floridi · lente de la infoesfera}
\textit{«Somos organismos informacionales ---inforgs--- que compartimos un entorno global hecho, en última instancia, de información: la infoesfera.»}

{\footnotesize\itshape\color{milcNegro!70}--- Luciano Floridi · La cuarta revolución (2014)}

Antes los mandatos llegaban de la casa, la escuela y el barrio ---unas pocas fuentes, y contradictorias entre sí, lo cual dejaba rendijas---. \textbf{Hoy llegan además de un entorno que no descansa y que está calibrado:} \emph{lo que se le muestra a alguien depende de lo que ya miró}, así que \textbf{a quien empieza a ver contenido sobre cómo debe ser un hombre o una mujer se le muestra más de lo mismo, cada vez más extremo}. \emph{Es la lógica del feed que trabajaste en el momento 7 del grado 8, aplicada aquí.} \textbf{Y la salida es la misma y es concreta:} \emph{el algoritmo aprende de lo que miras, así que lo que dejas de mirar también enseña.}

\textbf{De lo que me aparece en redes sobre «cómo debe ser» alguien de mi género, ¿qué he mirado para que me lo muestren?}
\end{softbox}
\linead{\linewidth}\\[1mm]
\linead{\linewidth}

\begin{infoband}{milcNegro}
\textbf{Nota para el docente.} Las tres son citas textuales y llevan su referencia debajo. Puedes remitir a la obra en clase.
\end{infoband}

\milcestacion{milcVino}{Mi autoevaluación · marca una casilla por fila}

{\small
\setlength{\extrarowheight}{2.2pt}
\begin{tabularx}{\textwidth}{>{\RaggedRight\arraybackslash\bfseries}p{3.0cm}YYY}
\rowcolor{milcVino}\color{white}\bfseries Criterio & \color{white}\bfseries Lo logré & \color{white}\bfseries En proceso & \color{white}\bfseries Necesito apoyo\\[.6mm]
Comprensión del tema & \milcopcion{Explico las ideas con mis palabras.} & \milcopcion{Reconozco las ideas, falta precisión.} & \milcopcion{Necesito repasar lo básico.}\\[.8mm]
\rowcolor{milcGris}Calidad del inventario & \milcopcion{Hice las dos listas, conté las horas y escogí un mandato mío para desobedecer con una acción.} & \milcopcion{Reconozco los mandatos ajenos, pero no encuentro ninguno que me apriete a mí.} & \milcopcion{Sigo pensando que eso ya cambió y hoy todos son libres de hacer lo que quieran.}\\[.8mm]
Aplicación y producto & \milcopcion{Mi producto cumple los criterios.} & \milcopcion{Está armado, con ajustes pendientes.} & \milcopcion{Aún está en construcción.}\\[.8mm]
\rowcolor{milcGris}Reflexión 5 dimensiones & \milcopcion{Respondí cada una con honestidad.} & \milcopcion{Respondí algunas con detalle.} & \milcopcion{Mis respuestas son muy generales.}\\[.8mm]
Triángulo de pensamiento & \milcopcion{Conecté las 3 voces con mi trabajo.} & \milcopcion{Conecté al menos una.} & \milcopcion{No las relaciono con mi proceso.}\\[.8mm]
\end{tabularx}}

\vspace{3mm}
\textbf{Hoy entendí que:}\\[1mm]
\linead{\linewidth}\\[1mm]
\linead{\linewidth}

\textbf{Mi compromiso para la próxima sesión es:}\\[1mm]
\linead{\linewidth}\\[1mm]
\linead{\linewidth}

\begin{softbox}{milcMostaza}{milcGris}{Cita de despedida · Marco Aurelio}
\textit{``El obstáculo en el camino se vuelve el camino. Lo que estorba al camino se vuelve camino.''}\\[1mm]
Cada error de hoy, bien observado, se convierte en pieza del camino que pisarás mañana.
\end{softbox}

\small
\begin{tabularx}{\textwidth}{>{\bfseries}p{3.6cm}Y}
\rowcolor{milcGradeDark!12}Nombre completo: & \linead{10cm}\\
Fecha: & \linead{4cm} \quad Curso/grupo: \linead{4cm}\\
Mi firma: & \linead{9cm}\\
\end{tabularx}
\normalsize

\begin{infoband}{milcGradeDark}
\textbf{Para la próxima sesión.} Dos cosas. \textbf{Primero, trae la tabla de horas sumada} y \textbf{el resultado de tu desobediencia} ---qué hiciste, qué pasó, quién dijo algo---. \emph{Y si no te atreviste, tráelo escrito también: qué te lo impidió es información valiosa.} \textbf{Segundo, pregunta en tu casa cómo era el reparto:} averigua con \textbf{una mujer mayor de tu familia} \textbf{qué le dejaron estudiar o hacer y qué no}, \textbf{quién lo decidió} y si alguna vez quiso algo que no le permitieron; y con \textbf{un hombre mayor}, \textbf{qué no podía hacer o decir}, y si alguna vez le tocó callarse algo que le dolía. \textbf{Trae:} tu tabla y las dos respuestas. \emph{Vas a encontrar dos historias que casi nunca se cuentan juntas: casi todas las familias tienen una mujer que quiso estudiar algo y no la dejaron, y casi todas tienen un hombre que nunca pudo decir que estaba mal. Las dos son el mismo mandato.}
\end{infoband}

\end{document}
